% GENERATED FILE. Edit canonical lessons, then run npm run generate:books.
% canonical-source-hash: fnv1a64:7d876d12414ebfb0
% canonical-lessons: RU-C51-cup, RU-C51-plate, RU-C51-spoon, RU-C51-knife, RU-C51-fork

\chapter{Cup, Plate, Spoon, Knife, Fork}
\label{ch:ru-cup-plate-spoon-knife-fork}
\hlchaptermodality{51}

\begin{chapteropening}
\textbf{By the end of this chapter:} \emph{I can name a cup, a plate, a spoon, a knife and a fork.}

Complete the last lesson of chapter 51: I can name a cup, a plate, a spoon, a knife and a fork.
\end{chapteropening}

\section[\texorpdfstring{\ru{чашка} (cháshka)}{чашка (cháshka)}]{\ru{чашка} (cháshka) — a cup}

\label{lesson:RU-C51-cup}



\begin{quote}
Before the new one: say the Russian for the bathroom, then the Russian for a clock.
\end{quote}

\subsection*{You'll want to know: \ru{чашка}}
\textbf{\ru{чашка}} (\emph{cháshka}) — \textquotedblleft{}a cup\textquotedblright{}. For \textbf{\ru{чай}} or \textbf{\ru{кофе}}.

\begin{cousinweb}
\textbf{\ru{чашка}} is \emph{\ru{чаша}}, \textquotedblleft{}a bowl, a goblet\textquotedblright{}, with \emph{-\ru{ка}}, \textquotedblleft{}little\textquotedblright{}. \emph{\ru{чаша}} $\leftarrow$ Proto-Slavic \emph{*čaša}, which Vasmer thinks is probably related to Old Prussian \emph{kiosi}, \textquotedblleft{}a goblet\textquotedblright{}.
\end{cousinweb}

\begin{grammarlens}[title={What you've built}]
One more word for this chapter.
\end{grammarlens}

\subsection*{Your turn}
\begin{itemize}
\raggedright
  \item \emph{Say it:} \emph{cháshka}
  \item \emph{Say it:} \emph{cháshka}, once more
  \item \emph{Say it:} say \emph{cháshka}
  \item \emph{Recall:} say the Russian for the bathroom, then the Russian for a clock, then say \emph{cháshka} again
\end{itemize}

\begin{tcolorbox}[breakable,colback=teal!4,colframe=teal!35!black,title={Before you move on}]
What does \ru{чашка} mean? (\textquotedblleft{}A cup\textquotedblright{}.) Say it once more.
\end{tcolorbox}
\section[\texorpdfstring{\ru{тарелка} (tarélka)}{тарелка (tarélka)}]{\ru{тарелка} (tarélka) — a plate}

\label{lesson:RU-C51-plate}



\begin{quote}
Before the new one: say the Russian for a clock, then the Russian for a cup.
\end{quote}

\subsection*{You'll want to know: \ru{тарелка}}
\textbf{\ru{тарелка}} (\emph{tarélka}) — \textquotedblleft{}a plate\textquotedblright{}.

\begin{cousinweb}
\textbf{\ru{тарелка}} — older \emph{\ru{тарель}}, recorded from 1509; Vasmer thinks it apparently came through Polish \emph{talerz} and Czech \emph{talíř} from Middle High German \emph{talier}, from Italian \emph{tagliere}, \textquotedblleft{}a trencher\textquotedblright{}, from \emph{tagliare}, \textquotedblleft{}to cut\textquotedblright{}.
\end{cousinweb}

\begin{grammarlens}[title={What you've built}]
One more word for this chapter.
\end{grammarlens}

\subsection*{Your turn}
\begin{itemize}
\raggedright
  \item \emph{Say it:} \emph{tarélka}
  \item \emph{Say it:} \emph{tarélka}, once more
  \item \emph{Say it:} say \emph{tarélka}
  \item \emph{Recall:} say the Russian for a clock, then the Russian for a cup, then say \emph{tarélka} again
\end{itemize}

\begin{tcolorbox}[breakable,colback=teal!4,colframe=teal!35!black,title={Before you move on}]
What does \ru{тарелка} mean? (\textquotedblleft{}A plate\textquotedblright{}.) Say it once more.
\end{tcolorbox}
\section[\texorpdfstring{\ru{ложка} (lózhka)}{ложка (lózhka)}]{\ru{ложка} (lózhka) — a spoon}

\label{lesson:RU-C51-spoon}



\begin{quote}
Before the new one: say the Russian for a cup, then the Russian for a plate.
\end{quote}

\subsection*{You'll want to know: \ru{ложка}}
\textbf{\ru{ложка}} (\emph{lózhka}) — \textquotedblleft{}a spoon\textquotedblright{}.

\begin{cousinweb}
\textbf{\ru{ложка}} $\leftarrow$ Old Russian \emph{lŭžĭka}, beside \emph{lŭžica}, from Proto-Slavic \emph{*lŭžĭka} — inherited. Vasmer thinks it is probably related to Albanian \emph{lugë}, \textquotedblleft{}spoon\textquotedblright{}, and perhaps to a root \textquotedblleft{}to break, to split\textquotedblright{}, since Swedish \emph{sked}, \textquotedblleft{}spoon\textquotedblright{}, goes with a verb \textquotedblleft{}to split\textquotedblright{}.
\end{cousinweb}

\begin{grammarlens}[title={What you've built}]
One more word for this chapter.
\end{grammarlens}

\subsection*{Your turn}
\begin{itemize}
\raggedright
  \item \emph{Say it:} \emph{lózhka}
  \item \emph{Say it:} \emph{lózhka}, once more
  \item \emph{Say it:} say \emph{lózhka}
  \item \emph{Recall:} say the Russian for a cup, then the Russian for a plate, then say \emph{lózhka} again
\end{itemize}

\begin{tcolorbox}[breakable,colback=teal!4,colframe=teal!35!black,title={Before you move on}]
What does \ru{ложка} mean? (\textquotedblleft{}A spoon\textquotedblright{}.) Say it once more.
\end{tcolorbox}
\section[\texorpdfstring{\ru{нож} (nozh)}{нож (nozh)}]{\ru{нож} (nozh) — a knife}

\label{lesson:RU-C51-knife}



\begin{quote}
Before the new one: say the Russian for a plate, then the Russian for a spoon.
\end{quote}

\subsection*{You'll want to know: \ru{нож}}
\textbf{\ru{нож}} (\emph{nozh}) — \textquotedblleft{}a knife\textquotedblright{}.

\begin{cousinweb}
\textbf{\ru{нож}} $\leftarrow$ Proto-Slavic \emph{*nožĭ} — inherited. Vasmer will not separate it from \emph{\ru{вонзить}}, \textquotedblleft{}to thrust in\textquotedblright{}, and \emph{\ru{заноза}}, \textquotedblleft{}a splinter\textquotedblright{}, and reports the usual link with Greek \emph{nýssō}, \textquotedblleft{}I prick\textquotedblright{}: the knife as the thing that pierces.
\end{cousinweb}

\begin{grammarlens}[title={What you've built}]
One more word for this chapter.
\end{grammarlens}

\subsection*{Your turn}
\begin{itemize}
\raggedright
  \item \emph{Say it:} \emph{nozh}
  \item \emph{Say it:} \emph{nozh}, once more
  \item \emph{Say it:} say \emph{nozh}
  \item \emph{Recall:} say the Russian for a plate, then the Russian for a spoon, then say \emph{nozh} again
\end{itemize}

\begin{tcolorbox}[breakable,colback=teal!4,colframe=teal!35!black,title={Before you move on}]
What does \ru{нож} mean? (\textquotedblleft{}A knife\textquotedblright{}.) Say it once more.
\end{tcolorbox}
\section[\texorpdfstring{\ru{вилка} (vílka)}{вилка (vílka)}]{\ru{вилка} (vílka) — a fork}

\label{lesson:RU-C51-fork}



\begin{quote}
Before the new one: say the Russian for a spoon, then the Russian for a knife.
\end{quote}

\subsection*{You'll want to know: \ru{вилка}}
\textbf{\ru{вилка}} (\emph{vílka}) — \textquotedblleft{}a fork\textquotedblright{}.

\begin{cousinweb}
\textbf{\ru{вилка}} is \emph{\ru{вилы}}, \textquotedblleft{}a pitchfork\textquotedblright{}, with \emph{-\ru{ка}}, \textquotedblleft{}little\textquotedblright{}: a little pitchfork. \emph{\ru{вилы}} $\leftarrow$ Proto-Slavic \emph{*vidla}, first \textquotedblleft{}something bent\textquotedblright{}, from the root of \emph{\ru{вить}}, \textquotedblleft{}to twist\textquotedblright{} (Vasmer).
\end{cousinweb}

\begin{grammarlens}[title={What you've built}]
One more word for this chapter.
\end{grammarlens}

\subsection*{Your turn}
\begin{itemize}
\raggedright
  \item \emph{Say it:} \emph{vílka}
  \item \emph{Say it:} \emph{vílka}, once more
  \item \emph{Say it:} say \emph{vílka}
  \item \emph{Recall:} say the Russian for a spoon, then the Russian for a knife, then say \emph{vílka} again
\end{itemize}

\begin{tcolorbox}[breakable,colback=teal!4,colframe=teal!35!black,title={Before you move on}]
What does \ru{вилка} mean? (\textquotedblleft{}A fork\textquotedblright{}.) Say it once more.
\end{tcolorbox}
